\documentclass[10pt,a4paper]{amsart}
\usepackage[margin=19mm]{geometry}
\usepackage{amsmath,amssymb,mathtools}
\usepackage[scr=rsfs,cal=euler]{mathalfa}
\usepackage{xfrac}
\usepackage[unicode,hidelinks,bookmarks=false]{hyperref}
\newcommand{\ie}{i.e.\ }
\newcommand{\cf}{cf.\ }
\newcommand{\resp}{respectively\ }
\newcommand{\Subsection}{\subsection}
\providecommand{\nobreakdash}{\nobreak\hbox{-}}
\setlength{\emergencystretch}{2em}
\multlinegap0pt
\usepackage{bm}
\usepackage{bbm}
\usepackage{dsfont}
\usepackage{xspace}
\usepackage{cancel}
\usepackage{mathtools}
\usepackage{cleveref}
\usepackage{amsthm}
\makeatletter
\@ifundefined{thm}{\newtheorem{thm}{Thm}}{}
\@ifundefined{lemma}{\newtheorem{lemma}[thm]{Lemma}}{}
\@ifundefined{proposition}{\newtheorem{proposition}[thm]{Proposition}}{}
\makeatother
\newcommand{\muminus}{\mu_{-}}
\newcommand{\rect}{\mathrm{rec}}
\title[A superpotential for Grassmannian Schubert varieties]{A superpotential for Grassmannian Schubert varieties}
\author{Konstanze Rietsch, Lauren Williams}
\date{Source: arXiv:2409.00734v2 (CC BY 4.0)}
\begin{document}
\maketitle

\noindent\emph{Source and licence.} A superpotential for Grassmannian Schubert varieties. Version arXiv:2409.00734v2 (CC BY 4.0), distributed under the Creative Commons Attribution 4.0 International licence, as recorded in the arXiv licence metadata for this exact version and shown on the abstract page https://arxiv.org/abs/2409.00734v2. Full rights notice: licenses/arxiv-2409.00734-CC-BY-4.0.txt.\par\smallskip

\noindent\textit{Editorial scope.} Counted unit: the Proposition whose source label is \texttt{p:GammaShiftRect} in the article (source0.tex lines 5657--5661, source label \texttt{p:GammaShiftRect}), delivered together with the article's own complete original proof of it (source0.tex lines 5663--5720, one \texttt{proof} environment of 58 lines). No mathematics is written, repaired or re-derived: statement and proof are the article's own text line for line. Required context retained uncounted: l:rectangles (lines 3839--3856); eq:1gen (lines 3843--3855); eq:2gen (lines 3843--3855); eq:3gen (lines 3843--3855); eq:4gen (lines 3843--3855); e:rr'Cartier (lines 5640--5642). Every definition, display and label of those lines is retained unchanged. Omitted from this package and not redistributed: 1--3838, 3856--5639, 5643--5656, 5662--5662, 5721--8605. Nothing inside a retained excerpt is omitted or altered: every retained body is delivered exactly as the article's lines are written, so no omission receipt of any kind accompanies this package. Citations: the retained excerpts carry no citation command at all, so no bibliography entry of the article is transcribed and no reference is redistributed. Reference adaptation: none was needed and none was made; every reference inside the delivered unit resolves inside it, so no unresolved reference or citation is produced. Printed slips kept literal: no printed word, symbol, hypothesis or number of the counted statement or of its proof was altered. Layout and class: the article's own class is \texttt{amsart} with packages including quiver, subfig, adjustbox, enumerate, etex, pstricks, amssymb, pict2e, inputenc, tikz, tikz-cd, geometry, rotating, graphicx; the wrapper is the lane's pinned \texttt{amsart} editorial wrapper at 10pt on A4, which restates the theorem-like environments the retained text needs and adds the article's own macro definitions that the retained text needs (\texttt{\textbackslash muminus}, \texttt{\textbackslash rect}), with extra wrapper packages (bm, bbm, dsfont, xspace, cancel, mathtools, cleveref). The delivered numbering is the unit's own. Assets: no retained span contains a figure environment, a graphics-inclusion command, a PDF-page inclusion, a table or a TikZ picture; no image file is copied into this package, the package declares no asset dependency and no image file is redistributed with it. Licence: version arXiv:2409.00734v2 of this article is distributed under the Creative Commons Attribution 4.0 International licence, as recorded in the arXiv licence metadata for this exact version and shown on the abstract page https://arxiv.org/abs/2409.00734v2. A full rights notice is delivered at licenses/arxiv-2409.00734-CC-BY-4.0.txt. Characters: every retained excerpt is pure ASCII, so the delivered engine, XeLaTeX with its default Latin Modern text font, reports no missing character.
\begin{lemma}\label{l:rectangles}
	Let  $\mathbf{r}=(r_1,\dots,r_d)$ and $\mathbf{r'}=(r'_1,\dots,r'_{n-1})$. Recall that $\lambda$ has $n-k$ rows and $k$ columns. 
	The superpotential polytope $\Gamma^{\lambda}_{\rect}(\mathbf{r},\mathbf{r'})$ 
is cut out by the following inequalities:
	\begin{align}
		v_{\mu} - v_{\muminus} & \leq r_\ell\ \text{ whenever }\mu 
		\text{ labels the $\ell$th outer corner of }\lambda  \label{eq:1gen}\\
		0 & \leq v_{1 \times 1}+r'_{n-k} \label{eq:2gen} \\
		v_{i\times j} - v_{(i-1)\times (j-1)} &\leq
		v_{(i+1) \times j}  - v_{i\times (j-1)}+r'_{n-k-i}
\ \text{ for }1 \leq i, 
		1\leq j, \text{ and } ((i+1) \times j) \subseteq \lambda \label{eq:3gen}\\
v_{i\times j}-v_{(i-1)\times (j-1)} &\leq
		v_{i\times (j+1)}  - v_{(i-1)\times j} + r'_{n-k+j}
\ \text{ for } 1\leq i,
		1 \leq j, \text{ and } (i \times (j+1)) \subseteq \lambda. \label{eq:4gen}
	\end{align}
\end{lemma}

\begin{equation}\label{e:rr'Cartier}
r_j=R-\sum_{b'_{i}\in NW(b_{\rho_{2j-1}})}r'_i
\end{equation}

\begin{proposition}\label{p:GammaShiftRect} Suppose that $R>0$ and $(\mathbf r,\mathbf r')$ satisfy \eqref{e:rr'Cartier}. We have that
\[
\Gamma^\lambda_{\rect}(\mathbf r,\mathbf r')=R\Gamma^\lambda_\rect(\mathbf 1,\mathbf 0)-v_\rect(\mathbf r').
\]
\end{proposition}

\begin{proof} 
Let $G=G_{\lambda}^{\rect}$ for the duration of this proof. 
 \cref{l:rectangles} gives us the following explicit description of $R\Gamma^\lambda_G(\mathbf 1,\mathbf 0)$.
\begin{align}
		v_{\mu} - v_{\muminus} & \le R\quad \text{ whenever }\mu 
		\text{ labels an outer corner of }\lambda  \label{eq:1}\\
		0 & \leq v_{1 \times 1} \label{eq:2} \\
		v_{(i-1)\times j} - v_{(i-2)\times (j-1)} &\leq
v_{i \times j}  - v_{(i-1)\times (j-1)}
\ \text{ for }2 \leq i, 
		1\leq j, \text{ and } (i \times j) \subseteq \lambda \label{eq:3}\\	
v_{i\times (j-1)}-v_{(i-1)\times (j-2)} &\leq
v_{i\times j}  - v_{(i-1)\times (j-1)}
\ \text{ for } 1\leq i,
		2 \leq j, \text{ and } (i \times j) \subseteq \lambda. \label{eq:4}
	\end{align}
	
	The polytope $R\Gamma^\lambda_G(\mathbf 1,\mathbf 0)-v_G(\mathbf r')$ can be described by shifting each coordinate $v_\nu$ of a point in $R\Gamma^\lambda_G(\mathbf 1,\mathbf 0)$ to $v'_\nu=v_\nu- r'_{(\nu)}$. 

	Consider first a frozen index $\mu\in\mathcal P_G$ associated to the outer corner box $b_{\rho_{2\ell-1}}$. Using the definition of the $r'_{(\nu)}$ and the inequality \eqref{eq:1} we get 
\[
v'_\mu-v'_{\muminus}=v_\mu-v_{\muminus}-r'_\mu+r'_{\muminus}=v_\mu-v_{\muminus}-\sum_{b'_{i}\in NW(b_{\rho_{2\ell-1}})}r'_i\le R -\sum_{b'_{i}\in NW(b_{\rho_{2\ell-1}})}r'_i,
\]
Therefore, 
 using \eqref{e:rr'Cartier}, the equivalent inequality to \eqref{eq:1} for the translated polytope becomes
\begin{equation} v'_{\mu} - v'_{\muminus} \le r_\ell, \label{eq:1'}
 \end{equation}
 which agrees with the inequality \eqref{eq:1gen} for $\Gamma^\lambda_G(\mathbf r,\mathbf{r'})$.
  
Since $v'_{1\times 1}=v_{1\times 1}-r'_{n-k}$ we get the inequality
\begin{equation}-r'_{n-k}\le  v'_{1\times 1}  \label{eq:2'}
 \end{equation}
 for $v'_{1\times 1}$, equivalent to \eqref{eq:2gen}.
 
We then have
\begin{equation*}
\begin{array}{ll}
v'_{i \times j}  - v'_{(i-1)\times (j-1)}&=  v_{i \times j}  - v_{(i-1)\times (j-1)} - r'_{i\times j}+r'_{(i-1)\times (j-1)}\\
&=v_{i \times j}  - v_{(i-1)\times (j-1)} -r'_{n-k-i+1}-r'_{n-k-i+2}+\dotsc- r'_{n-k+j-1}\\
v'_{(i-1) \times j}  - v'_{(i-2)\times (j-1)}&=  v_{(i-1) \times j}  - v_{(i-2)\times (j-1)}- r'_{(i-1)\times j}+r'_{(i-2)\times (j-1)}\\
&=v_{(i-1) \times j}  - v_{(i-2)\times (j-1)}- r'_{n-k-i+2}-r'_{n-k-i+3}-\dotsc - r'_{n-k+j-1}
	\end{array}
	\end{equation*}
	and therefore 
	\begin{multline*}
	\left(v'_{i \times j}  - v'_{(i-1)\times (j-1)}\right)-\left(v'_{(i-1) \times j}  - v'_{(i-2)\times (j-1)}\right)=\\
	\left(v_{i \times j}  - v_{(i-1)\times (j-1)}\right)-\left(v_{(i-1) \times j}  - v_{(i-2)\times (j-1)}\right)-r'_{n-k-i+1}.
	\end{multline*}
	Therefore the equivalent inequality to \eqref{eq:3} for the translated polytope is 
	\begin{equation}
	v'_{(i-1)\times j} - v'_{(i-2)\times (j-1)} \leq
v'_{i \times j}  - v'_{(i-1)\times (j-1)}+r'_{n-k-i+1}
\ \text{ for }2 \leq i, 
		1\leq j, \text{ and } (i \times j) \subseteq \lambda, \label{eq:3'}
		\end{equation}
		which agrees with the inequality~\eqref{eq:3gen} for $\Gamma^\lambda_G(\mathbf r,\mathbf r')$. 
		The completely analogous calculation shows the inequality \eqref{eq:4gen} is the shifted version of \eqref{eq:4}. Thus we have shown the statement of the lemma.
\end{proof}

\end{document}
